% Cuentas Nacionales de Transferencia y envejecimiento · Parte 2: la economía generacional
% Finanzas Públicas, Maestría en Economía Aplicada, Tec de Monterrey
% Esquema aprobado: esquema_nta.md (2026-11-01). Compilar con: latexmk -pdf nta_p2_manual.tex
% Versión para repartir: nta_p2_manual_handout.tex (misma fuente, opción handout).
% Las notas del presentador van en \note{}; para verlas, compilar con \setbeameroption{show notes}.
% Datos de las gráficas: datos/prep_p2.py.
\documentclass[aspectratio=169,11pt]{beamer}

\usepackage[T1]{fontenc}
\usepackage{lmodern}
\usepackage[spanish,es-nodecimaldot]{babel}
\unaccentedoperators % la opción es-noaccentedoperators no existe en esta versión de babel
\usepackage{amsmath,amssymb}
\usepackage{booktabs}
\usepackage{array}
\usepackage{tikz}
\usetikzlibrary{babel,arrows.meta,calc,positioning,decorations.pathreplacing,patterns}
\usepackage{pgfplots}
\pgfplotsset{compat=1.18}
\usepgfplotslibrary{fillbetween}

% ---------------------------------------------------------------------------
% Tema sobrio (igual al de las sesiones anteriores)
% ---------------------------------------------------------------------------
\definecolor{cTitulo}{RGB}{0,55,100}
\definecolor{cOferta}{RGB}{31,87,153}
\definecolor{cDemanda}{RGB}{200,100,0}
\definecolor{cImpuesto}{RGB}{175,30,30}
\definecolor{cRecauda}{RGB}{40,130,80}
\definecolor{cGris}{RGB}{110,110,110}

\usetheme{default}
\setbeamertemplate{navigation symbols}{}
\setbeamercolor{frametitle}{fg=cTitulo}
\setbeamercolor{title}{fg=cTitulo}
\setbeamercolor{structure}{fg=cTitulo}
\setbeamercolor{block title}{bg=cTitulo!12,fg=cTitulo}
\setbeamercolor{block body}{bg=cTitulo!4}
\setbeamercolor{block title alerted}{bg=cImpuesto!12,fg=cImpuesto}
\setbeamercolor{block body alerted}{bg=cImpuesto!4}
\setbeamerfont{frametitle}{size=\large,series=\bfseries}
\setbeamertemplate{itemize items}[circle]
\setbeamertemplate{blocks}[rounded]
\setbeamertemplate{footline}{%
  \hfill{\small\color{cGris}NTA y envejecimiento · Parte 2\quad\insertframenumber/\inserttotalframenumber}\hspace*{1em}\vspace{1.5mm}}

% Fuente al pie de cada diapositiva
\newcommand{\fuente}[1]{\par\vfill{\small\color{cGris}#1\par}}
% Distintivo de diapositiva opcional
\newcommand{\opcional}{\tikz[remember picture,overlay]\node[anchor=north east,font=\small,text=cGris] at ([xshift=-5mm,yshift=-3mm]current page.north east) {opcional};}

% ---------------------------------------------------------------------------
% Estilo común de gráficas
% ---------------------------------------------------------------------------
\tikzset{
  every node/.append style={font=\small},
  eje/.style={-{Stealth[length=2mm]},thick,black!75},
  flecha/.style={-{Stealth[length=2mm]},thick,cTitulo},
  guia/.style={thin,dashed,black!45},
  sincorte/.style={execute at begin node={\hyphenpenalty=10000\exhyphenpenalty=10000\relax}},
  caja/.style={sincorte,draw=cTitulo,rounded corners,fill=cTitulo!5,align=center,inner sep=4pt},
  cajafisco/.style={sincorte,draw=cRecauda,rounded corners,fill=cRecauda!10,align=center,inner sep=3pt},
}
\pgfplotsset{
  estilo/.style={axis lines=left,tick label style={font=\small},label style={font=\small},
    legend style={font=\small,draw=none,fill=none},every axis plot/.append style={thick}},
}

\title{Cuentas Nacionales de Transferencia y envejecimiento}
\subtitle{Parte 2: la economía generacional}
\author{Héctor Juan Villarreal}
\institute{Finanzas Públicas · Maestría en Economía Aplicada\\Escuela de Gobierno y Transformación Pública · Tec de Monterrey}
\date{Noviembre de 2026}

\begin{document}

% =========================================================================
% I. LA ECONOMÍA GENERACIONAL · 22 min · cap. 1
% =========================================================================
\section{La economía generacional}

% 1 -----------------------------------------------------------------------
\begin{frame}[plain]
  \titlepage
  \note{1 min. Lectura: Naciones Unidas (2013), \emph{National Transfer Accounts Manual}, caps. 1 y 2. Existe edición en español: Naciones Unidas (2022), CEPAL.}
\end{frame}

% 2 -----------------------------------------------------------------------
\begin{frame}{Consumimos toda la vida y producimos solo una parte}
  \centering
  \begin{tikzpicture}
    \begin{axis}[estilo,name=a,width=6.8cm,height=4.0cm,xmin=0,xmax=90,ymin=0,ymax=1.8,xtick={0,20,40,60,80},
        title={Filipinas, 1999},title style={font=\small\bfseries},xlabel={Edad},
        legend style={at={(0.99,0.99)},anchor=north east},legend cell align=left]
      \addplot[cDemanda] table[col sep=comma,x=edad,y=C]{datos/p2_fil1999.csv};
      \addlegendentry{Consumo}
      \addplot[cOferta] table[col sep=comma,x=edad,y=YL]{datos/p2_fil1999.csv};
      \addlegendentry{Ingreso laboral}
    \end{axis}
    \begin{axis}[estilo,at={(a.east)},anchor=west,xshift=0.7cm,width=6.8cm,height=4.0cm,xmin=0,xmax=90,ymin=0,ymax=1.8,xtick={0,20,40,60,80},
        title={Alemania, 2003},title style={font=\small\bfseries},xlabel={Edad}]
      \addplot[cDemanda] table[col sep=comma,x=edad,y=C]{datos/p2_ale2003.csv};
      \addplot[cOferta] table[col sep=comma,x=edad,y=YL]{datos/p2_ale2003.csv};
    \end{axis}
  \end{tikzpicture}
  \par\smallskip\raggedright
  \begin{itemize}
    \item De niños y de viejos consumimos más de lo que ganamos con nuestro trabajo. En medio, al revés.
    \pause
    \item Filipinas: déficit antes de los 27 y desde los 60. Alemania: antes de los 27 y desde los 58.
  \end{itemize}
  \fuente{Perfiles por persona de ntaccounts.org; ingreso laboral medio de 30--49 años = 1. Manual NTA, pp.~3--4, fig.~1.1.}
  \note{2 min. Redibujada con los datos de la base: datos/p2\_fil1999.csv, datos/p2\_ale2003.csv. La fig. 1.1 del Manual usa los mismos países y años. Los cruces de edad salen iguales con estos datos.}
\end{frame}

% 3 -----------------------------------------------------------------------
\begin{frame}{La economía generacional: quién produce, quién consume y quién comparte}
  \begin{columns}[T]
    \begin{column}{0.45\textwidth}
      \begin{block}{Qué es}
        Las instituciones y los mecanismos con los que cada edad produce, consume, comparte y ahorra, y los flujos que resultan entre edades.
      \end{block}
      \pause
      \begin{alertblock}{Lo que no ven las cuentas nacionales}
        Dicen cuánto consume el país. No dicen a qué edad.
      \end{alertblock}
    \end{column}
    \begin{column}{0.52\textwidth}
      \pause
      \textbf{Cinco razones para medirlo}
      \begin{enumerate}
        \item Niños y muchos mayores dependen de otros para no caer en pobreza.
        \item La edad ya no significa lo mismo: se estudia más y a los 65 se está más sano.
        \item El crecimiento depende de la escuela y la salud de los niños.
        \item Los sistemas de apoyo se presionan si la estructura por edad cambia.
        \item Surge la pregunta de equidad: ¿les estamos cargando de más a los jóvenes?
      \end{enumerate}
    \end{column}
  \end{columns}
  \fuente{Manual NTA, pp.~1--3.}
  \note{2 min. La definición es de Mason y Lee (2011), citada en la p.~3.}
\end{frame}

% 4 -----------------------------------------------------------------------
\begin{frame}{Tres formas de cubrir el déficit}
  \centering
  \begin{tabular}{@{}l>{\raggedright\arraybackslash}p{5.2cm}>{\raggedright\arraybackslash}p{5.2cm}@{}}
    \toprule
    & \textbf{Transferencias} & \textbf{Activos} \\
    & {\small sin contraprestación} & {\small hoy se guarda, mañana se usa} \\
    \midrule
    \textbf{Público} & La escuela, el hospital, la pensión de reparto & Deuda pública, fondos soberanos \\[0.3em]
    \textbf{Privado} & Lo que los padres gastan en sus hijos, el dinero que mandan los hijos, las remesas & La casa propia, los ahorros para el retiro, el crédito \\
    \bottomrule
  \end{tabular}
  \par\medskip\raggedright
  \pause
  \begin{itemize}
    \item La familia mueve recursos por costumbre y cariño. El Estado, por ley.
    \item Con activos no se le pide nada a nadie: se compra oro de joven y se vende de viejo. O se pide prestado de joven y se paga después.
  \end{itemize}
  \pause
  \begin{block}{Entonces}
    Tres vías: transferencias públicas, transferencias privadas y activos.
  \end{block}
  \fuente{Manual NTA, pp.~3--5, cuadro~1.1.}
  \note{3 min. El Manual cuenta cuatro casillas con dos mecanismos y dos sectores; en la práctica se reportan tres vías (fig. 1.6). El capital público aparece como casi nulo porque las cuentas nacionales valúan su producto al costo (p.~5).}
\end{frame}

% 5 -----------------------------------------------------------------------
\begin{frame}{Cada país paga la vejez de otra manera}
  \begin{columns}[T]
    \begin{column}{0.56\textwidth}
      \textbf{Déficit de los de 65 y más: quién lo cubre (\%)}\par\smallskip
      \begin{tabular}{@{}lrrr@{}}
        \toprule
        & Estado & Familia & Activos \\
        \midrule
        México, 2004 & 37 & $-$25 & \textbf{89} \\
        Japón, 2004 & 57 & 1 & 42 \\
        Alemania, 2003 & 72 & $-$7 & 36 \\
        Costa Rica, 2013 & 74 & $-$10 & 36 \\
        Suecia, 2005 & \textbf{115} & $-$14 & $-$1 \\
        \bottomrule
      \end{tabular}
      \par\medskip
      {\small Un número negativo es recurso que sale: en México, los mayores daban a su familia más de lo que recibían.}
    \end{column}
    \begin{column}{0.41\textwidth}
      \begin{itemize}
        \item México: la casa y lo ahorrado.
        \item Suecia: el Estado, todo y algo más.
        \item La familia casi nunca sostiene en neto; Taiwán, Tailandia y Corea son la excepción.
      \end{itemize}
      \pause
      \begin{alertblock}{Pregunta}
        Nuestra cuenta 2016--2024 da 54.5\,\% de Estado para los de 60+. ¿Qué cambió: México o la medición?
      \end{alertblock}
    \end{column}
  \end{columns}
  \fuente{Cálculo propio con ntaccounts.org y ONU, WPP 2024. Manual NTA, pp.~15--16.}
  \note{3 min. No se pueden redibujar los puntos del triángulo de la fig. 1.6 porque el Manual no da valores. Calculé con la base: públicas, privadas y reasignaciones por activos entre déficit, 65+, con pesos WPP (datos/prep\_p2.py). Coincide con el Manual: México con dos tercios o más en activos y familia negativa. Respuesta posible: las dos. Cambió el programa no contributivo y cambió el método; además aquí son 60+, no 65+.}
\end{frame}

% 6 -----------------------------------------------------------------------
\begin{frame}{La razón de sostenimiento: cuántos producen por cada uno que consume}
  \begin{columns}[T]
    \begin{column}{0.53\textwidth}
      \[
        \frac{L}{N}=\frac{\sum_a y_a\,P_a}{\sum_a c_a\,P_a}
        \quad
        \begin{array}{l}
          \text{productores efectivos}\\
          \text{consumidores efectivos}
        \end{array}
      \]
      $P_a$: personas de edad $a$. $y_a$ y $c_a$: ingreso laboral y consumo de esa edad, con 30--49 años $=1$.
      \par\medskip
      \pause
      \begin{block}{Ejemplo: China, 2025}
        1.40 mil millones de consumidores efectivos y 0.73 de productores efectivos: $L/N=0.52$. Un productor por cada dos consumidores.
      \end{block}
    \end{column}
    \begin{column}{0.44\textwidth}
      \pause
      \textbf{Por qué no contar personas de 15 a 64}
      \begin{itemize}
        \item Un joven de 16 que estudia no produce como uno de 45.
        \item Una persona de 80 no consume como un niño.
        \item Sirve para evaluar políticas: si la gente se retira más tarde, sube $L$ y sube la razón.
      \end{itemize}
    \end{column}
  \end{columns}
  \fuente{Manual NTA, pp.~10--11; cap.~5, cuadro~5.6. En inglés, \emph{support ratio}; el Manual en español dice «razón de sustento económico». Usamos el término de CEPAL.}
  \note{3 min. La p.~11 del Manual dice que el retiro tardío «reduce» la razón: es errata, la aumenta. Los pesos se fijan en un año base. El ejemplo de China está en el cap. 5, fuera de la lectura; lo uso porque el cap. 1 solo da la idea en palabras.}
\end{frame}

% 7 -----------------------------------------------------------------------
\begin{frame}{Los dos bonos, en una identidad}
  \begin{columns}[T]
    \begin{column}{0.54\textwidth}
      Consumo por consumidor:
      \[
        \frac{C}{N}=\frac{(1-s)\,Y}{L}\cdot\frac{L}{N}
      \]
      \pause
      En tasas de crecimiento:
      \[
        gr\!\left[\frac{C}{N}\right]=
        \underbrace{gr\!\left[\frac{(1-s)Y}{L}\right]}_{\text{segundo bono}}+
        \underbrace{gr\!\left[\frac{L}{N}\right]}_{\text{primer bono}}
      \]
      \pause
      {\small Si $L/N$ sube 1 punto, el consumo por consumidor sube 1 punto, dado lo demás.}
    \end{column}
    \begin{column}{0.43\textwidth}
      \pause
      \textbf{Primer bono.} China, 1972--2014: la razón de sostenimiento sumó 0.7\,\% al año. De 2014 a 2050 restará 0.5\,\% al año.
      \par\medskip
      \pause
      \textbf{Segundo bono.} Depende de qué se haga con el primero:
      \begin{itemize}
        \item \emph{Capital humano:} menos hijos, más gasto por hijo.
        \item \emph{Capital:} más ahorro para una vejez larga.
      \end{itemize}
      \par\smallskip
      No tiene fecha ni es automático.
    \end{column}
  \end{columns}
  \fuente{Manual NTA, pp.~10--14, ecs.~1.1--1.2.}
  \note{3 min. Es lo que en la Parte 1 vimos en intuición. El gasto en capital humano por hijo va de 100--200\,\% del ingreso laboral de 30--49 años en países de alta fecundidad a más de 500\,\% en algunos de baja (pp.~12--13, fig. 1.4). Elasticidad un poco menor que 1.}
\end{frame}

% 8 -----------------------------------------------------------------------
\begin{frame}{Las flechas de Lee: hacia dónde va cada peso público}
  \opcional
  \begin{columns}[T]
    \begin{column}{0.56\textwidth}
      \centering
      \begin{tikzpicture}[x=0.105cm,y=0.00125cm]
        \draw[eje] (10,0) -- (72,0);
        \foreach \v in {10,20,30,40,50,60,70} {\draw (\v,0) -- (\v,-50) node[below] {\v};}
        \node at (41,-500) {Edad};
        \foreach \yb/\h/\a/\b/\c in {0/1138/38.1/29.1/cGris,1138/762/38.0/62.5/cImpuesto,1900/288/38.9/33.2/cRecauda,2188/294/37.6/16.0/cOferta} {
          \pgfmathsetmacro\s{ifthenelse(\b>\a,1,-1)}
          \pgfmathsetmacro\tp{\b-\s*2.4}
          \pgfmathsetmacro\ym{\yb+\h/2}
          \pgfmathsetmacro\yt{\yb+\h-25}
          \fill[\c!70] (\a,\yb) -- (\tp,\yb) -- (\b,\ym) -- (\tp,\yt) -- (\a,\yt) -- cycle;
        }
        \node[anchor=west,font=\small] at (39,2335) {Educación};
        \node[anchor=west,font=\small] at (39.5,2045) {Salud};
        \node[anchor=west,font=\small,text=white] at (39,1519) {Pensiones};
        \node[anchor=west,font=\small] at (39,569) {Otros};
      \end{tikzpicture}
      \par\smallskip
      {\small\color{cGris}Altura: flujo anual por persona. La flecha va de la edad media de quien paga a la de quien recibe.}
    \end{column}
    \begin{column}{0.41\textwidth}
      \textbf{Brasil, 2002}\par\smallskip
      \begin{tabular}{@{}lrr@{}}
        \toprule
        & De & A \\
        \midrule
        Educación & 37.6 & 16.0 \\
        Salud & 38.9 & 33.2 \\
        Pensiones & 38.0 & 62.5 \\
        Otros & 38.1 & 29.1 \\
        \bottomrule
      \end{tabular}
      \pause
      \begin{block}{El área importa}
        Altura por distancia es riqueza de transferencias: 18,642 reales por persona en pensiones. Alguien que todavía no nace la está debiendo.
      \end{block}
    \end{column}
  \end{columns}
  \fuente{Manual NTA, pp.~17--18, fig.~1.8, redibujada con sus valores.}
  \note{4 min, opcional. Valores: educación 294, salud 288, pensiones 762, otros 1,138 reales al año. Riqueza: $-$6,361; $-$1,638; $+$18,642; $-$10,277; total 365 reales, unos 4\,\% del ingreso de un adulto de 30--49 (p.~18). Las flechas hacia abajo dan riqueza negativa: casi todos ya recibieron la escuela y ahora la pagan. Es pariente de las Cuentas Generacionales de Auerbach y Kotlikoff.}
\end{frame}

% 9 -----------------------------------------------------------------------
\begin{frame}{La razón de sostenimiento fiscal: la cuenta que le espera a Hacienda}
  \begin{columns}[T]
    \begin{column}{0.50\textwidth}
      \[
        \frac{\text{contribuyentes efectivos}}{\text{beneficiarios efectivos}}
      \]
      Cada edad se pondera con lo que paga y con lo que recibe, comparada con los de 30--49.
      \par\medskip
      \centering
      \begin{tikzpicture}[x=1cm,y=1.55cm]
        \fill[cOferta!70] (0,0) rectangle (1.2,1);
        \fill[cImpuesto!70] (2,0) rectangle (3.2,0.7);
        \draw[eje] (-0.2,0) -- (3.7,0);
        \node[above] at (0.6,1) {1.0};
        \node[above] at (2.6,0.7) {cerca de 0.7};
        \node[below] at (0.6,0) {2010};
        \node[below] at (2.6,0) {2050};
      \end{tikzpicture}
      \par{\small\color{cGris}Brasil, con 2010 = 1.}
    \end{column}
    \begin{column}{0.47\textwidth}
      \begin{itemize}
        \item Brasil, 2010 a 2050: solo por la estructura por edad, hay que bajar beneficios 30\,\% o subir impuestos otro tanto.
      \end{itemize}
      \pause
      \begin{block}{¿Y si lo pagamos con deuda?}
        No resuelve nada: agrega intereses, y esos también se pagan con impuestos.
      \end{block}
      \pause
      \begin{alertblock}{Pregunta}
        Si les toca en Hacienda: ¿menos pensión, menos salud o más impuestos?
      \end{alertblock}
    \end{column}
  \end{columns}
  \fuente{Manual NTA, pp.~18--19, fig.~1.9.}
  \note{2 min. El Manual da el 30\,\%, no el valor exacto de la serie; por eso la barra dice «cerca de 0.7». Brasil tuvo la mayor caída de 20 economías (Miller, 2011). Tailandia no cae porque casi no tiene programas públicos para mayores: la razón depende del diseño.}
\end{frame}

% =========================================================================
% II. EL PANORAMA DE LAS CUENTAS · 33 min · cap. 2
% =========================================================================
\section{El panorama de las cuentas}

% 10 ----------------------------------------------------------------------
\begin{frame}{Qué miden las NTA, y qué no}
  \begin{columns}[T]
    \begin{column}{0.49\textwidth}
      \textbf{Qué son}
      \begin{itemize}
        \item Todos los flujos económicos del año, repartidos por edad.
        \item La unidad es la persona. Empresas y gobierno son intermediarios.
        \item Los totales cuadran con las cuentas nacionales.
      \end{itemize}
      \pause
      \begin{alertblock}{Lo que no son}
        Son contabilidad de un año, de corte transversal. No dicen qué causa qué, ni siguen a una generación en el tiempo.
      \end{alertblock}
    \end{column}
    \begin{column}{0.48\textwidth}
      \pause
      \textbf{Cuatro diferencias con el SCN}
      \begin{enumerate}
        \item Manda la edad; las instituciones son secundarias.
        \item El ingreso mixto se reparte: dos tercios al trabajo, un tercio al capital.
        \item Todo se mide antes de impuestos a los productos.
        \item Se estiman los flujos dentro del hogar.
      \end{enumerate}
    \end{column}
  \end{columns}
  \fuente{Manual NTA, pp.~6, 25--29, 34--35.}
  \note{3 min. Por eso «ingreso laboral» no existe tal cual en el SCN: ni el ingreso mixto ni los impuestos a los productos se tratan igual. Se usa el SCN como control de los totales (recuadro 2.2).}
\end{frame}

% 11 ----------------------------------------------------------------------
\begin{frame}{La identidad de flujos: lo que entra es igual a lo que sale}
  Para el grupo de edad $x$, en el año:
  \[
    \underbrace{Y^l(x)+\tau^+(x)+Y^k(x)+Y^{p+}(x)}_{\text{entradas}}
    =
    \underbrace{C(x)+\tau^-(x)+Y^{p-}(x)+S(x)}_{\text{salidas}}
  \]
  \pause
  Pasamos consumo e ingreso laboral de un lado, y lo demás del otro:
  \[
    \underbrace{C(x)-Y^l(x)}_{\text{déficit del ciclo de vida}}
    =
    \underbrace{\tau^+(x)-\tau^-(x)}_{\text{transferencias netas}}
    +
    \underbrace{Y^A(x)-S(x)}_{\text{reasignaciones por activos}}
  \]
  \pause
  con $Y^A=Y^k+Y^{p+}-Y^{p-}$: capital más ingreso neto de propiedad.
  \pause
  \begin{block}{Cómo se lee}
    Lo que falta (o sobra) a esa edad se cubre con transferencias y activos. Es una identidad, no una teoría: no dice que una cosa cause la otra.
  \end{block}
  \fuente{Manual NTA, p.~29, ecs.~2.1--2.2. El ahorro $S$ es la partida que cierra.}
  \note{4 min. Vale por persona, en agregado y por edad. Transferencias y rentas de la propiedad tienen contraparte: suman cero con el resto del mundo (pp.~30--31).}
\end{frame}

% 12 ----------------------------------------------------------------------
\begin{frame}{El ciclo de vida económico, en dólares}
  \begin{columns}[T]
    \begin{column}{0.60\textwidth}
      \textbf{Estados Unidos, 2003, dólares por persona}\par\smallskip
      \begin{tabular}{@{}lrrr@{}}
        \toprule
        & 0--24 & 25--64 & 65+ \\
        \midrule
        \textbf{Consumo} & \textbf{23,951} & \textbf{34,069} & \textbf{43,734} \\
        \quad Público & 9,451 & 5,771 & 14,310 \\
        \quad\quad Educación & 4,842 & 194 & 0 \\
        \quad\quad Salud & 699 & 1,666 & 10,400 \\
        \quad Privado & 14,500 & 28,298 & 29,424 \\
        \quad\quad Educación & 913 & 344 & 76 \\
        \quad\quad Salud & 1,322 & 3,524 & 5,417 \\
        \textbf{Ingreso laboral} & \textbf{3,425} & \textbf{43,233} & \textbf{5,689} \\
        \quad Salarios & 3,265 & 37,390 & 4,287 \\
        \quad Autoempleo & 160 & 5,842 & 1,402 \\
        \bottomrule
      \end{tabular}
    \end{column}
    \begin{column}{0.37\textwidth}
      \begin{itemize}
        \item Los niños consumen escuela.
        \item Los mayores, salud: 15,817 dólares, más de un tercio de su consumo.
      \end{itemize}
      \pause
      \begin{block}{El consumo público cuenta}
        La escuela gratuita y la consulta del IMSS son consumo de quien las usa. Las paga otro: de eso sale la transferencia.
      \end{block}
    \end{column}
  \end{columns}
  \fuente{Manual NTA, pp.~35 y 39, cuadro~2.4.}
  \note{3 min. 10,400 + 5,417 = 15,817; entre 43,734 da 36\,\%. El consumo público se reparte por edad; la parte «otros» del consumo público es igual para todas las edades (3,910).}
\end{frame}

% 13 ----------------------------------------------------------------------
\begin{frame}{El déficit del ciclo de vida, y una sorpresa mexicana}
  \begin{columns}[T]
    \begin{column}{0.50\textwidth}
      \textbf{Estados Unidos, 2003, dólares}\par\smallskip
      \begin{tabular}{@{}lrrr@{}}
        \toprule
        & 0--24 & 25--64 & 65+ \\
        \midrule
        Por persona & 20,526 & $-$9,164 & 38,045 \\
        Billones & 2.10 & $-$1.65 & 1.30 \\
        \bottomrule
      \end{tabular}
      \par\medskip
      \begin{itemize}
        \item Por persona pesan más los mayores. En total, los jóvenes: son el triple.
        \item Suma positiva: el ingreso laboral no alcanza. El resto sale de los activos.
      \end{itemize}
    \end{column}
    \begin{column}{0.47\textwidth}
      \pause
      \centering
      \begin{tikzpicture}
        \begin{axis}[estilo,width=6.3cm,height=3.7cm,xmin=0,xmax=90,ymin=0,ymax=1.4,xtick={0,20,40,60,80},
            xlabel={Edad},title={México, 2014},title style={font=\small\bfseries},
            legend style={at={(0.5,1.02)},anchor=north,legend columns=2},legend cell align=left]
          \addplot[cDemanda] table[col sep=comma,x=edad,y=C]{datos/p2_mex2014.csv};
          \addlegendentry{Consumo}
          \addplot[cOferta] table[col sep=comma,x=edad,y=YL]{datos/p2_mex2014.csv};
          \addlegendentry{Ingreso laboral}
        \end{axis}
      \end{tikzpicture}
      \par\raggedright
      El promedio mexicano solo tiene superávit a los 45 y 46 años (en 2004, de los 33 a los 48). Nuestra cuenta da déficit en todas las edades: al trabajo le toca poco del PIB.
    \end{column}
  \end{columns}
  \fuente{Manual NTA, pp.~33, 38; cuadros~2.2--2.3. México: ntaccounts.org; López Reynosa y Villarreal (2026), trabajo en curso.}
  \note{4 min. El superávit en México es un hallazgo de esta preparación: datos/p2\_mex2014.csv y datos/p2\_mex2004.csv (ingreso medio 30--49 = 1). 2010: 37 a 50 años. Es la serie de Mejía Guevara en ntaccounts.org, independiente de nuestra base. En nuestra cuenta: el consumo incluye el del gobierno y el ingreso laboral es 39\,\% a 44\,\% del PIB; el perfil masculino sí tiene superávit, de 33 a 62 años (resumen\_ejecutivo.tex). No es que los mexicanos consuman de más.}
\end{frame}

% 14 ----------------------------------------------------------------------
\begin{frame}{Reasignaciones públicas: la pensión y la consulta del abuelo}
  \centering
  \begin{tikzpicture}
    \begin{axis}[estilo,name=a,width=6.8cm,height=3.5cm,xmin=0,xmax=90,ymin=0,ymax=1.3,stack plots=y,area style,xtick={0,20,40,60,80},
        title={Brasil, 2008},title style={font=\small\bfseries}]
      \addplot[fill=cOferta!60,draw=none] table[col sep=comma,x=edad,y=edu]{datos/p2_publicas_bra2008.csv} \closedcycle;
      \addplot[fill=cRecauda!60,draw=none] table[col sep=comma,x=edad,y=sal]{datos/p2_publicas_bra2008.csv} \closedcycle;
      \addplot[fill=cImpuesto!60,draw=none] table[col sep=comma,x=edad,y=pen]{datos/p2_publicas_bra2008.csv} \closedcycle;
      \addplot[fill=cGris!40,draw=none] table[col sep=comma,x=edad,y=otros]{datos/p2_publicas_bra2008.csv} \closedcycle;
    \end{axis}
    \begin{axis}[estilo,at={(a.east)},anchor=west,xshift=0.7cm,width=6.8cm,height=3.5cm,xmin=0,xmax=90,ymin=0,ymax=1.3,stack plots=y,area style,xtick={0,20,40,60,80},
        title={México, 2004},title style={font=\small\bfseries}]
      \addplot[fill=cOferta!60,draw=none] table[col sep=comma,x=edad,y=edu]{datos/p2_publicas_mex2004.csv} \closedcycle;
      \addplot[fill=cRecauda!60,draw=none] table[col sep=comma,x=edad,y=sal]{datos/p2_publicas_mex2004.csv} \closedcycle;
      \addplot[fill=cImpuesto!60,draw=none] table[col sep=comma,x=edad,y=pen]{datos/p2_publicas_mex2004.csv} \closedcycle;
      \addplot[fill=cGris!40,draw=none] table[col sep=comma,x=edad,y=otros]{datos/p2_publicas_mex2004.csv} \closedcycle;
    \end{axis}
  \end{tikzpicture}
  \par\vspace{-1mm}
  {\small
  \tikz\fill[cOferta!60] (0,0) rectangle (0.3,0.2); Educación\quad
  \tikz\fill[cRecauda!60] (0,0) rectangle (0.3,0.2); Salud\quad
  \tikz\fill[cImpuesto!60] (0,0) rectangle (0.3,0.2); Pensiones\quad
  \tikz\fill[cGris!40] (0,0) rectangle (0.3,0.2); Otros\quad
  {\color{cGris}(eje horizontal: edad)}}
  \par\smallskip\raggedright
  Estados Unidos, 2003: un mayor recibe 25,273 dólares al año, 80\,\% en salud y pensiones. Y paga 9,830.
  \pause
  \begin{alertblock}{Pregunta}
    El abuelo paga IVA en el restaurante y predial por su casa. ¿Qué parte de su pensión se paga él mismo?
  \end{alertblock}
  \fuente{Lo que recibe cada persona; ingreso laboral de 30--49 = 1. ntaccounts.org; Manual NTA, pp.~16--17, 40.}
  \note{4 min. Base del bloque central de la Parte 3. A los 75 años, en Brasil salud y pensiones son 0.99 de 1.12; en México 2004, 0.17 de 0.35 (la serie mexicana trae 0.10 más en otras transferencias en efectivo). Esto es 2004: todavía no existía la pensión universal actual. Datos: datos/p2\_publicas\_bra2008.csv y p2\_publicas\_mex2004.csv. 10,400 + 9,835 = 20,235; entre 25,273 da 80.1\,\%. Cambié el ejemplo del esquema: las medicinas tienen tasa cero de IVA.}
\end{frame}

% 15 ----------------------------------------------------------------------
\begin{frame}{Reasignaciones privadas: lo que pasa dentro de la casa}
  \begin{columns}[T]
    \begin{column}{0.56\textwidth}
      \textbf{Transferencias netas, EUA, 2003, dólares}\par\smallskip
      \begin{tabular}{@{}lrrr@{}}
        \toprule
        & 0--24 & 25--64 & 65+ \\
        \midrule
        Dentro del hogar & 11,088 & $-$6,479 & $-$2,324 \\
        Entre hogares & 127 & $-$254 & $-$909 \\
        \midrule
        \textbf{Total privado} & \textbf{11,215} & \textbf{$-$6,733} & \textbf{$-$3,233} \\
        \bottomrule
      \end{tabular}
      \par\medskip
      \begin{itemize}
        \item Los padres sostienen a los hijos. Casi todo ocurre dentro de la casa, y ninguna encuesta lo pregunta.
        \item Los mayores también dan: el flujo neto sale de ellos.
      \end{itemize}
    \end{column}
    \begin{column}{0.41\textwidth}
      \pause
      \begin{block}{Cómo se estima lo que nadie reporta}
        \begin{itemize}
          \item Se reparte el consumo del hogar con una regla: los adultos consumen igual; los niños, menos.
          \item El jefe del hogar recibe todo el ingreso de activos y el ahorro.
        \end{itemize}
      \end{block}
      \pause
      Son supuestos fuertes: el jefe importa, y suele ser hombre.
    \end{column}
  \end{columns}
  \fuente{Manual NTA, pp.~27--28, 41 (cuadro~2.8) y 53.}
  \note{3 min. Salud y educación se tratan aparte en la regla de reparto (p.~27). En países con mucha corresidencia, el supuesto del jefe pesa más (p.~28).}
\end{frame}

% 16 ----------------------------------------------------------------------
\begin{frame}{Reasignaciones basadas en activos: ingreso de activos menos ahorro}
  \begin{columns}[T]
    \begin{column}{0.52\textwidth}
      \[
        \text{reasignación por activos}=Y^A-S
      \]
      \textbf{Privadas, Estados Unidos, 2003, 65+}\par\smallskip
      \begin{tabular}{@{}lr@{}}
        \toprule
        Ingreso de activos & 23,749 \\
        Ahorro (se resta) & $-$1,396 \\
        \midrule
        \textbf{Reasignación} & \textbf{25,146} \\
        \bottomrule
      \end{tabular}
      \par\smallskip
      {\small El ahorro es negativo: los mayores se comen lo guardado.}
    \end{column}
    \begin{column}{0.45\textwidth}
      \pause
      \begin{block}{Cierra con el déficit de 38,045}
        \begin{tabular}{@{}lr@{}}
          Transf.\ públicas & 15,443 \\
          Transf.\ privadas & $-$3,233 \\
          Activos públicos & 690 \\
          Activos privados & 25,146 \\
          \midrule
          \textbf{Suma} & \textbf{38,046} \\
        \end{tabular}
      \end{block}
      \pause
      Un detalle fiscal: los intereses de la deuda pública los registra el Manual como transferencia que sale de los contribuyentes.
    \end{column}
  \end{columns}
  \fuente{Manual NTA, pp.~5, 31, 38, 41; cuadros~2.2, 2.6 y 2.8. Dólares por persona; la diferencia de 1 es redondeo.}
  \note{3 min. Los flujos públicos por activos son el préstamo del gobierno: en el cuadro estilizado 2.1, el ahorro público es negativo. Los jóvenes de 25--34 ahorran más de lo que reciben de sus activos; en Corea la reasignación se vuelve positiva a fines de los 30 (fig. 1.2).}
\end{frame}

% 17 ----------------------------------------------------------------------
\begin{frame}{Población por perfiles: de cada persona a la presión agregada}
  \begin{columns}[T]
    \begin{column}{0.50\textwidth}
      \[
        \text{agregado}(x)=\text{perfil}(x)\times\text{población}(x)
      \]
      \begin{itemize}
        \item En Filipinas, el déficit de los niños es casi 15 veces el de los mayores.
        \item En Alemania, el de los mayores es 1.5 veces el de los niños.
      \end{itemize}
      \pause
      Gran parte de esa diferencia es la pirámide, no el comportamiento.
    \end{column}
    \begin{column}{0.47\textwidth}
      \pause
      \begin{block}{Las transferencias tienen que cuadrar}
        Lo que uno recibe, otro lo paga. Si cambia la estructura por edad, algún perfil se mueve.
      \end{block}
      \begin{itemize}
        \item Con menos niños por adulto: más por niño o menos carga para el adulto.
        \item Con más mayores por adulto: más carga, menos por mayor o ambas.
      \end{itemize}
    \end{column}
  \end{columns}
  \fuente{Manual NTA, pp.~3, 30 y 50--51.}
  \note{3 min. Por eso una proyección que multiplica perfiles fijos por la población futura es una primera aproximación: la restricción obliga a que algo se ajuste (pp.~50--51). Esa es la pregunta de la Parte 3.}
\end{frame}

% 18 ----------------------------------------------------------------------
\begin{frame}{Comparaciones: México se sale del patrón}
  \opcional
  \begin{columns}[T]
    \begin{column}{0.58\textwidth}
      \centering
      \begin{tikzpicture}
        \begin{axis}[estilo,width=7.2cm,height=5.0cm,xmin=0,xmax=90,ymin=-0.6,ymax=1.1,xtick={0,20,40,60,80},
            xlabel={Edad},ylabel={Consumo menos ingreso laboral},
            legend style={at={(0.5,1.02)},anchor=south,legend columns=4},legend cell align=left]
          \draw[guia] (axis cs:0,0) -- (axis cs:90,0);
          \addplot[cImpuesto,very thick] table[col sep=comma,x=edad,y=mex]{datos/p2_deficit_paises.csv}; \addlegendentry{México}
          \addplot[cOferta] table[col sep=comma,x=edad,y=bra]{datos/p2_deficit_paises.csv}; \addlegendentry{Brasil}
          \addplot[cRecauda] table[col sep=comma,x=edad,y=chl]{datos/p2_deficit_paises.csv}; \addlegendentry{Chile}
          \addplot[cDemanda] table[col sep=comma,x=edad,y=ale]{datos/p2_deficit_paises.csv}; \addlegendentry{Alemania}
        \end{axis}
      \end{tikzpicture}
    \end{column}
    \begin{column}{0.39\textwidth}
      \textbf{Superávit máximo}\par\smallskip
      \begin{tabular}{@{}lr@{}}
        \toprule
        Alemania, 2008 & 0.48 \\
        Chile, 2018 & 0.39 \\
        Brasil, 2008 & 0.23 \\
        México, 2014 & 0.01 \\
        \bottomrule
      \end{tabular}
      \par\medskip
      \pause
      Normalizar es lo que permite comparar: pesos, reales y euros se vuelven «veces el ingreso de un adulto de 30--49».
    \end{column}
  \end{columns}
  \fuente{Déficit por persona; ingreso laboral de 30--49 años = 1. ntaccounts.org, consultado el 1 de noviembre de 2026; Lee y Mason (2011).}
  \note{3 min, opcional. El cap. 2 no trae comparaciones normalizadas: por eso uso la base. Datos: datos/p2\_deficit\_paises.csv. Años distintos: lo honesto es decirlo. Otros superávits máximos: Costa Rica 2013, 0.31; Colombia 2014, 0.19; Perú 2014, 0.18; Argentina 2016, 0.33; Uruguay 2013, 0.25.}
\end{frame}

% 19 ----------------------------------------------------------------------
\begin{frame}{Lo que las NTA no ven}
  \begin{columns}[T]
    \begin{column}{0.49\textwidth}
      \begin{block}{El trabajo en casa}
        La cuenta básica no mide la producción del hogar. Cuidar a un niño o a un padre no aparece, y lo hacen sobre todo mujeres.
      \end{block}
      \begin{block}{El sexo}
        El método básico no separa hombres y mujeres, y sin el trabajo en casa una cuenta por sexo subestima a las mujeres.
      \end{block}
    \end{column}
    \begin{column}{0.48\textwidth}
      \begin{block}{La desigualdad dentro de cada edad}
        Da promedios. El «mayor promedio» no existe: la riqueza se concentra y la media queda arriba de la mediana.
      \end{block}
      \pause
      \begin{alertblock}{Por eso}
        Nuestro 76.3\,\% del déficit en mujeres mide el ciclo monetario, no lo que aportan. La medición del tiempo regresa en la Parte 3.
      \end{alertblock}
    \end{column}
  \end{columns}
  \fuente{Manual NTA, pp.~52--53; López Reynosa y Villarreal (2026), trabajo en curso.}
  \note{3 min. El Manual remite a Donehower (2011) para las cuentas con tiempo; hoy existe el manual de las Naciones Unidas de 2024.}
\end{frame}

% 20 ----------------------------------------------------------------------
\begin{frame}{Qué datos se necesitan}
  \opcional
  \begin{columns}[T]
    \begin{column}{0.49\textwidth}
      \textbf{Tres tipos de datos}
      \begin{enumerate}
        \item Población por edad.
        \item Cuentas nacionales: los totales.
        \item Encuestas y registros: la forma por edad.
      \end{enumerate}
      \pause
      \begin{block}{Controles macro}
        Si la encuesta da 100 y las cuentas nacionales 110, se multiplica cada edad por 1.1.
      \end{block}
    \end{column}
    \begin{column}{0.48\textwidth}
      \pause
      \textbf{En México}
      \begin{itemize}
        \item ENIGH: forma por edad.
        \item Cuentas Nacionales de INEGI: nivel.
      \end{itemize}
      \par\smallskip
      La encuesta capta entre 10\,\% y 53\,\% del agregado, según la variable. Los factores de ajuste van de 1.9 a 9.9.
      \par\medskip
      Por eso decimos: la encuesta da la forma; las cuentas nacionales, el nivel.
    \end{column}
  \end{columns}
  \fuente{Manual NTA, p.~28 (recuadro~2.2) y p.~56; López Reynosa y Villarreal (2026), trabajo en curso.}
  \note{3 min, opcional. La p.~56 está en el cap. 3, fuera de la lectura. El esquema mencionaba la Cuenta Pública: nuestro trabajo no la usa como control, así que no aparece. Fuente: Missions/Funded/NTA\_Inequality/paper\_final.tex.}
\end{frame}

% =========================================================================
% III. CIERRE · 5 min
% =========================================================================
\section{Cierre}

% 21 ----------------------------------------------------------------------
\begin{frame}{El presupuesto es un sistema de reasignaciones por edad}
  \begin{itemize}
    \item \textbf{Cada impuesto y cada gasto mueve recursos entre edades.} Una reforma de pensiones no es solo de pensiones.
    \pause
    \item \textbf{Con Kaplow} evaluamos cada política junto con su financiamiento. Aquí le agregamos la pregunta de a qué edad paga cada quien, y qué generación carga con la factura.
    \pause
    \item \textbf{El 10 de noviembre} vemos la parte macro. La deuda también mueve recursos entre generaciones, y por eso va de la mano de estas cuentas.
  \end{itemize}
  \fuente{Elaboración propia, con Manual NTA, pp.~16--22.}
  \note{3 min. Amarre con la sesión de Kaplow del 26 y 27 de octubre y con la de deuda del 10 de noviembre.}
\end{frame}

% 22 ----------------------------------------------------------------------
\begin{frame}{Para la Parte 3}
  \begin{columns}[T]
    \begin{column}{0.48\textwidth}
      Ya tenemos las piezas:
      \begin{enumerate}
        \item Un déficit en la niñez y otro en la vejez.
        \item Tres vías para cubrirlos.
        \item Un Estado que paga, sobre todo, pensiones y salud a los mayores.
      \end{enumerate}
    \end{column}
    \begin{column}{0.48\textwidth}
      \pause
      \begin{alertblock}{Pregunta}
        ¿Qué le pasa a ese sistema cuando la población envejece rápido?
      \end{alertblock}
    \end{column}
  \end{columns}
  \fuente{Sigue la Parte 3: Cecchini et al. (2025), CEPAL.}
  \note{2 min.}
\end{frame}

\end{document}
